\subsection{Principal Findings}

The results answer the three research questions at different levels of evidence. For RQ1, real PX4 telemetry supported runtime anomaly detection and, conditional on an anomalous window having a specific label, four-class fault-domain diagnosis. The fixed chronological protocol yielded strong detector and Stage 2 results, whereas the lower leave-log-out detector AUPRC showed that this capability is sensitive to flight-log isolation. The high anomaly-only Stage 2 Macro-F1 should therefore be interpreted as evidence that the labeled fault domains are distinguishable after anomaly selection, not as the performance of a complete online diagnostic system.

For RQ2, the hierarchical SHAP analysis provided quantitatively testable diagnostic evidence rather than only qualitative feature-importance examples. Evidence mass was conserved across feature, channel, topic, and subsystem aggregation, and masking the validation-ranked SHAP features degraded diagnostic performance substantially more than masking equal-size random subsets. However, the document-derived mapping sensitivity materially reduced overall Consistency@1, left External Position near random, and placed Global Position below its class-specific random baseline. The explanation layer is consequently useful for organizing model-relevant telemetry evidence, but semantic agreement is mapping-dependent and cannot be treated as uniform explanation validation.

The majority-commit uORB graph extended this evidence into software-module suspect rankings and thereby partially addressed the architecture-aware part of RQ2. Module aggregation was restricted to the 1,393 anomalous test windows whose firmware matched the graph. The explicit-declaration audit supported the parser within its auditable source scope, but degree normalization changed the leading module order, and publisher/subscriber relations still encode possible evidence paths rather than fault causality. The resulting ranks are most appropriately used to prioritize inspection of modules associated with influential topics, not to identify a defective module without independent module-level ground truth.

For RQ3, strict isolation, external data, and controlled replay exposed complementary limits. Purged and leave-log-out results reduced the apparent robustness of the fixed chronological evaluation; ALFA supported only partial zero-shot transfer of the binary detector; and the controlled mutations were not detected by the unchanged real-flight detector. Taken together, these results support a constrained claim: HS-AeroTS forms an auditable chain from runtime anomaly scores to fault-domain and architecture-aware evidence, but the evaluated chain does not establish reliable end-to-end software fault localization.
